\documentclass{article}
\usepackage{amsmath, amssymb, amsthm}
\usepackage{hyperref}
\usepackage{geometry}
\geometry{margin=1in}

\title{Erd\H{o}s Problem \#864}
\date{Last updated: 2025-08-31}
\author{Source: erdosproblems.com}

\begin{document}
\maketitle

\noindent\textbf{Status:} OPEN \quad \textbf{No prize}

\noindent\textbf{Tags:} number theory, sidon sets, additive combinatorics

\medskip
\noindent\textbf{Problem Statement:}

Let $A\subseteq \{1,\ldots N\}$ be a set such that there exists at most one $n$ with more than one solution to $n=a+b$ (with $a\leq b\in A$). Estimate the maximal possible size of $\lvert A\rvert$ - in particular, is it true that\[\lvert A\rvert \leq (1+o(1))\frac{2}{\sqrt{3}}N^{1/2}?\]

\bigskip
\noindent\textbf{Remarks:}

A problem of Erd\H{o}s and Freud, who prove that\[\lvert A\rvert \geq (1+o(1))\frac{2}{\sqrt{3}}N^{1/2}.\]This is shown by taking a genuine Sidon set $B\subset [1,N/3]$ of size $\sim N^{1/2}/\sqrt{3}$ and taking the union with $\{N-b : b\in B\}$.

For the analogous question with $n=a-b$ they prove that $\lvert A\rvert\sim N^{1/2}$.

This is a weaker form of [840].

\bigskip
\noindent\small{Source: \url{https://www.erdosproblems.com/864}}
\end{document}
